\section{From pairs to groups}\label{sec:groups}

Who flies with whom, as a group? The contact links two pilots at a time: it flickers, it is not
transitive, and it changes during the afternoon. This section turns it into groups in four steps:
pair intervals, groups at each instant (configurations), chains of groups (episodes), and the
flights tied by them (dependency blocks).

\subsection{Pair intervals}\label{sec:pair-intervals}

The instantaneous contact flickers wherever a quantity sits close to its threshold (a distance
near $\varepsilon_h$, a height difference near $\varepsilon_z$, an angle near $\theta_{\max}$), at
the changes of state, where the heading condition switches on and off, and where the phase
labelling is wrong on a few fixes.

\begin{definition}[pair interval]\label{def:pair-interval}
For each pair consider the maximal runs of consecutive grid instants with $a_{ij}=1$. Two runs
separated by a gap shorter than $g$ are joined; the resulting runs shorter than $\tau^*$ are
dropped. Each remaining run is a half-open interval $[t_{\rm on}, t_{\rm off})$, and the set of
the intervals of the day is $\I$.
\end{definition}

Joining the short gaps plays the role of the robustness of Buchin et al.~\cite[§4]{buchin2015},
which ignores interruptions shorter than $\alpha$; the maximal segments of Li et
al.~\cite{li2013,ruzzo1999} are an alternative. The minimum duration $\tau^*$ keeps chance
contacts out, and the null model fixes it (Section~\ref{sec:calibration}).

\begin{figure}
\centering
\begin{tikzpicture}[x=1.0cm,y=1cm,
    row/.style={font=\footnotesize, anchor=east}]
  \foreach \y in {2,1,0} \fill[light] (0,\y-0.17) rectangle (12,\y+0.17);
  \node[row] at (-0.1,2) {$a_{ij}(t)$};
  \node[row] at (-0.1,1) {gaps $<g$ filled};
  \node[row] at (-0.1,0) {pieces $<\tau^*$ dropped};
  \foreach \a/\b in {0.4/2.0, 2.3/4.4, 4.8/5.2, 7.0/7.5, 9.0/10.9, 11.2/11.4} \fill[aqua] (\a,1.83) rectangle (\b,2.17);
  \foreach \a/\b in {0.4/5.2, 7.0/7.5, 9.0/11.4} \fill[aqua] (\a,0.83) rectangle (\b,1.17);
  \foreach \a/\b in {2.0/2.3, 4.4/4.8, 10.9/11.2} \fill[pattern=north east lines, pattern color=white] (\a,0.83) rectangle (\b,1.17);
  \foreach \a/\b in {0.4/5.2, 9.0/11.4} \fill[aqua] (\a,-0.17) rectangle (\b,0.17);
  \draw[ink2, line width=0.7pt] (7.0,-0.17) -- (7.5,0.17) (7.0,0.17) -- (7.5,-0.17);
  \foreach \x/\l in {0.4/t_{\rm on}, 5.2/t_{\rm off}, 9.0/t_{\rm on}, 11.4/t_{\rm off}} \node[lbl, below] at (\x,-0.2) {$\l$};
  \draw[<->, ink2] (5.2,1.45) -- node[above, lbl] {gap $>g$} (7.0,1.45);
  \draw[|-|, ink2] (0,-1.05) -- node[below, lbl] {$g$} (0.9,-1.05);
  \draw[|-|, ink2] (1.6,-1.05) -- node[below, lbl] {$\tau^*$} (3.1,-1.05);
\end{tikzpicture}
\caption{From the contact to the pair intervals. Top: the contact of one pair, with the flicker
of the instantaneous condition. Middle: gaps shorter than $g$ are filled (hatched). Bottom: runs
shorter than $\tau^*$ are dropped; two intervals remain.}
\label{fig:pair-intervals}
\end{figure}

\subsection{The problem of non-transitivity}\label{sec:transitivity}

If 1 flies with 8 early in the afternoon and 8 with 14 later, 1 and 14 may never have met: contact
is symmetric but not transitive. A partition into groups needs an \emph{equivalence relation}
$\sim$, reflexive ($i\sim i$), symmetric ($i\sim j\Rightarrow j\sim i$) and transitive ($i\sim j$
and $j\sim k\Rightarrow i\sim k$), so any definition of groups must give up something
(Figure~\ref{fig:transitivity}):
\begin{enumerate}[label=(\alph*)]
  \item direct contact between all members: the transitive closure puts 1 and 14 in one group
        although they never met;
  \item the partition: groups may overlap, here $\{1,8\}$ and $\{8,14\}$, as the maximal groups
        of Buchin et al.~\cite{buchin2015} do;
  \item the static picture: the groups are defined instant by instant.
\end{enumerate}
We take (c): at every instant the groups are the connected components of the graph of active
contacts, as the components of~\cite{buchin2015}, and their evolution is described without
giving a group a lasting identity.

\begin{figure}
\centering
\begin{tikzpicture}[x=1cm,y=1cm]
  \begin{scope}
    \draw[->, ink2] (0,0) -- (7.3,0) node[right, lbl, ink] {$t$};
    \foreach \x/\h in {0/13:00, 2.4/14:00, 4.8/15:00, 7.0/16:00} {
      \draw[ink2] (\x,-0.06) -- (\x,0.06); \node[lbl, below] at (\x,-0.06) {\h};
    }
    \fill[blueI] (0.2,0.85) rectangle (3.4,1.15);
    \node[lbl, anchor=east] at (0.1,1.0) {1--8};
    \fill[orangeJ] (4.0,0.35) rectangle (6.8,0.65);
    \node[lbl, anchor=west] at (6.9,0.5) {8--14};
    \node[panel] at (0,1.65) {pair intervals};
  \end{scope}
  \begin{scope}[shift={(8.6,0)}]
    \node[panel] at (0,1.65) {(a) transitive closure};
    \draw[ink2, rounded corners=8pt, fill=light!60] (0,-0.3) rectangle (2.7,1.25);
    \node[flight] at (0.45,0.5) {1}; \node[flight] at (1.35,0.5) {8}; \node[flight] at (2.25,0.5) {14};
    \node[note, anchor=north west, align=left] at (0,-0.45) {one group all afternoon:\\1 and 14 never met};
  \end{scope}
  \begin{scope}[shift={(0,-3.4)}]
    \node[panel] at (0,1.65) {(b) overlapping groups};
    \draw[blueI, line width=0.9pt] (0.95,0.5) ellipse (0.95 and 0.6);
    \draw[orangeJ, line width=0.9pt] (1.85,0.5) ellipse (0.95 and 0.6);
    \node[flight] at (0.45,0.5) {1}; \node[flight] at (1.4,0.5) {8}; \node[flight] at (2.35,0.5) {14};
    \node[note, anchor=north west, align=left] at (0,-0.4) {8 belongs to two groups};
  \end{scope}
  \begin{scope}[shift={(4.6,-3.4)}]
    \node[panel] at (0,1.65) {(c) instant by instant (our choice)};
    \node[lbl] at (0.9,1.15) {13:30};
    \draw[edge] (0.45,0.45) -- (1.35,0.45);
    \node[inE1] at (0.45,0.45) {1}; \node[inE1] at (1.35,0.45) {8}; \node[flight] at (0.9,-0.35) {14};
    \node[lbl] at (3.6,1.15) {15:30};
    \draw[edge] (3.15,0.45) -- (4.05,0.45);
    \node[flight] at (3.6,-0.35) {1}; \node[inE2] at (3.15,0.45) {8}; \node[inE2] at (4.05,0.45) {14};
    \node[note, anchor=west, align=left] at (4.6,0.3) {groups are the connected\\components at each instant};
  \end{scope}
\end{tikzpicture}
\caption{Non-transitivity: 1--8 early in the afternoon, 8--14 later, and 1 and 14 never meet.
The three ways out: (a) one group of all three, (b) two overlapping groups, (c) the groups of
each instant, our choice.}
\label{fig:transitivity}
\end{figure}

\subsection{Configurations}\label{sec:configurations}

The active graph at instant $t$ is
\begin{equation}\label{eq:graph}
  \begin{gathered}
  G(t)=\bigl(V(t),E(t)\bigr),\qquad V(t)=\{i:\ t\in T^\dep_i\},\\
  E(t)=\bigl\{\{i,j\}:\ \exists\,(i,j,t_{\rm on},t_{\rm off})\in\I,\ t\in[t_{\rm on},t_{\rm off})\bigr\}:
  \end{gathered}
\end{equation}
its nodes are the pilots that at $t$ have departed, have not landed and have a position; its edges
are the pairs whose interval contains $t$. A connected component of $G(t)$ with at least two nodes
is a group at instant $t$: its pilots are joined by chains of contacts (Appendix~\ref{app:graphs}
recalls the terms).

\begin{definition}[configuration]
A configuration $c$ is given by a set of flights $M_c$, its members, with $|M_c|\ge2$, and by an
interval of time $[t^{\rm start}_c, t^{\rm end}_c)$, such that $M_c$ is a connected component of
$G(t)$ at every $t$ of the interval and the interval is maximal. We write
$c=\bigl(M_c,[t^{\rm start}_c,t^{\rm end}_c)\bigr)$.
\end{definition}

The members of a configuration are fixed: a pilot who joins or leaves ends it and starts another.
Only the edges among the members may change: an edge that closes a triangle does not create a new
configuration (Figure~\ref{fig:configurations}). The configurations are the edges of the Reeb
graph of~\cite{buchin2015} with at least two entities (Figure~\ref{fig:reeb}), built on delayed
contacts instead of simultaneous proximity.

\begin{figure}
\centering
\begin{tikzpicture}[x=1cm,y=1cm]
  \foreach \sx/\t/\tri/\four in {0/t_1/0/0, 4.3/t_2/1/0, 8.6/t_3/1/1} {
    \begin{scope}[shift={(\sx,0)}]
      \node[lbl] at (1.3,2.55) {$G(\t)$};
      \coordinate (a) at (0.3,1.0); \coordinate (b) at (1.0,2.0); \coordinate (c) at (1.7,1.0); \coordinate (d) at (2.8,1.0);
      \draw[edge] (a) -- (b) -- (c);
      \ifnum\tri=1 \draw[edge] (a) -- (c); \fi
      \ifnum\four=1 \draw[edge] (c) -- (d); \fi
      \ifnum\four=1
        \node[inE2] at (a) {1}; \node[inE2] at (b) {2}; \node[inE2] at (c) {3}; \node[inE2] at (d) {4};
      \else
        \node[inE1] at (a) {1}; \node[inE1] at (b) {2}; \node[inE1] at (c) {3}; \node[flight] at (d) {4};
      \fi
    \end{scope}
  }
  \begin{scope}[shift={(0,-2.9)}]
    \draw[->, ink2] (0,0) -- (12.4,0) node[right, lbl, ink] {$t$};
    \foreach \x/\l in {1.3/t_1, 5.6/t_2, 9.9/t_3} {\draw[ink2, densely dotted] (\x,0) -- (\x,4.9); \node[lbl, below] at (\x,0) {$\l$};}
    \foreach \a/\b/\y/\n in {0.3/11.8/2.1/1--2, 0.3/11.8/1.7/2--3, 4.6/11.8/1.3/1--3, 8.6/11.8/0.9/3--4} {
      \fill[ink2] (\a,\y-0.09) rectangle (\b,\y+0.09);
      \node[lbl, anchor=east] at (0.2,\y) {\n};
    }
    \fill[blueI] (0.3,0.25) rectangle (8.55,0.6); \node[font=\footnotesize, text=white] at (4.4,0.425) {$c$: $\{1,2,3\}$};
    \fill[orangeJ] (8.65,0.25) rectangle (11.8,0.6); \node[font=\footnotesize, text=white] at (10.2,0.425) {$c'$: $\{1,2,3,4\}$};
    \node[lbl, anchor=east] at (0.2,0.425) {config.};
  \end{scope}
\end{tikzpicture}
\caption{Configurations. Top: the active graph at three instants. Bottom: the pair intervals of
the four flights and the configurations they produce. At $t_2$ the edge 1--3 closes a triangle:
the component keeps its members and $c$ goes on. When 4 joins, $c$ ends and $c'$ starts.}
\label{fig:configurations}
\end{figure}

\begin{proposition}
At every instant each flight belongs to at most one configuration.
\end{proposition}

\begin{proof}
The connected components of $G(t)$ form a partition of $V(t)$.
\end{proof}

The membership of flight $i$ at instant $t$ is therefore well defined:
\begin{equation}\label{eq:membership}
  \gamma_i(t)=\begin{cases}
    \st{L} & \text{if } k^0_i\le t<k^\dep_i \text{ (local phase), or the flight is local},\\
    c & \text{if } i\in M_c \text{ and } t\in[t^{\rm start}_c,t^{\rm end}_c),\\
    \st{S} & \text{if } t\in T^\dep_i \text{ otherwise (solo flight).}
  \end{cases}
\end{equation}
\fix{In the tables, every instant between the departure and the landing outside a configuration
is labelled \st{S}, the instants in a gap of the track included, where the flight has no position.
The analyses that read the fixes see nothing there; only the time fractions computed from the
membership include them.}

\subsection{Episodes}\label{sec:episodes}

A group that gains or loses a pilot becomes a new configuration, yet its pilots go on flying
together. Episodes keep this continuity.

\begin{definition}[episode]
Two configurations are linked, $c\to c'$, if one starts exactly when the other ends and they have
at least one member in common:
\begin{equation}\label{eq:link}
  c\to c' \iff t^{\rm start}_{c'}=t^{\rm end}_c\ \wedge\ M_c\cap M_{c'}\ne\emptyset .
\end{equation}
An episode is a connected component of the undirected graph associated with $\to$.
\end{definition}

So an episode is a chain of configurations without interruption: each starts when the previous one
ends and keeps at least one of its members, and two configurations of the same episode need not
share anyone. In the language of~\cite{buchin2015}, it is a connected component of the Reeb graph
without the edges of single pilots. Four properties matter (Figure~\ref{fig:episodes}):
\begin{enumerate}[label=(\roman*)]
  \item a split produces two configurations linked to the same previous one, so they stay in the
        same episode; the same holds for a merge;
  \item if the group dissolves completely, the chain breaks and the episode ends; a group that
        forms again later opens a new episode;
  \item no threshold is needed: one member in common is enough;
  \item the same flight can belong to several episodes of the same day, at different times.
\end{enumerate}

\begin{figure}
\centering
\begin{tikzpicture}[x=1.4cm,y=1cm,
    conf/.style={draw=none, minimum height=5mm, inner sep=0pt, font=\footnotesize, text=white},
    link/.style={->, ink, line width=0.7pt}]
  \draw[->, ink2] (0,-0.4) -- (9.6,-0.4) node[right, lbl, ink] {$t$};
  \def\confbox#1#2#3#4#5{\fill[#5] (#1,#3-0.25) rectangle (#2,#3+0.25); \node[font=\footnotesize, text=white] at ({(#1+#2)/2},#3) {#4};}
  \confbox{0}{1.95}{2.6}{$c_1\,\{1,2\}$}{blueI}
  \confbox{0}{1.95}{0.9}{$c_2\,\{3,4\}$}{blueI}
  \confbox{2.05}{3.95}{1.75}{$c_3\,\{1,2,3,4\}$}{blueI}
  \confbox{4.05}{5.95}{2.6}{$c_4\,\{1,2\}$}{blueI}
  \confbox{4.05}{5.45}{0.9}{$c_5\,\{3,4\}$}{blueI}
  \confbox{7.5}{9.0}{1.75}{$c_6\,\{2,5\}$}{orangeJ}
  \draw[link] (1.95,2.6) -- (2.05,1.85);
  \draw[link] (1.95,0.9) -- (2.05,1.65);
  \draw[link] (3.95,1.85) -- (4.05,2.6);
  \draw[link] (3.95,1.65) -- (4.05,0.9);
  \node[note] at (2.0,3.15) {merge};
  \node[note] at (4.0,3.15) {split};
  \node[note, align=center] at (6.6,1.25) {dissolved:\\the chain breaks};
  \node[note, align=center] at (8.25,0.85) {new episode;\\flight 2 again};
  \draw[decorate, decoration={brace, mirror, amplitude=5pt}, blueI] (0,0.2) -- (5.95,0.2) node[midway, below=5pt, lbl, ink] {episode E1};
  \draw[decorate, decoration={brace, mirror, amplitude=5pt}, orangeJ] (7.5,0.2) -- (9.0,0.2) node[midway, below=5pt, lbl, ink] {E2};
\end{tikzpicture}
\caption{Episodes. Arrows are the links~\eqref{eq:link}. The merge of $c_1$ and $c_2$ into $c_3$
and the split of $c_3$ into $c_4$ and $c_5$ keep all five configurations in episode E1, although
$c_1$ and $c_5$ share no member. After the group dissolves, $c_6$ has no predecessor and opens E2:
flight 2 belongs to both episodes.}
\label{fig:episodes}
\end{figure}

\begin{remark}[why not a persistent group name]
One could keep the same group name while the overlap of the members between two instants exceeds a
threshold, as methods that track communities in evolving networks do with the Jaccard
index~\cite{greene2010}. The choice has three defects: the threshold is arbitrary; a group whose
members are replaced one at a time keeps its name even when none of the first members is left;
in splits and merges no natural criterion says who inherits the name. So we let the episode
record only the continuity, and the configuration the composition.
\end{remark}

\subsection{Dependency blocks}\label{sec:blocks}

Which flights can the statistics treat as independent? Those that never flew in the same episode,
directly or through a chain of episodes.

\begin{definition}[block]
For each day consider the bipartite graph flights--episodes, with an edge between a flight and
every episode it belongs to. Its connected components, restricted to the flights, are the
\emph{dependency blocks}; a flight never in a group is a block on its own.
\end{definition}

\begin{figure}
\centering
\begin{tikzpicture}[x=1.35cm,y=1cm,
    ep/.style={rectangle, draw=ink, line width=0.6pt, minimum size=5.6mm, inner sep=1pt, font=\footnotesize, fill=white}]
  \foreach \f/\x/\s in {1/0/inE1, 2/1/inE1, 3/2/inE1, 4/3/inE1, 5/4/flight, 6/5/inE2, 7/6/inE2, 8/7/inE2}
    \node[\s] (f\f) at (\x,1.6) {\f};
  \foreach \e/\x in {1/1, 2/2.5, 3/5.5, 4/6.5} \node[ep] (e\e) at (\x,0) {E\e};
  \foreach \e/\f in {1/1, 1/2, 1/3, 2/3, 2/4, 3/6, 3/7, 4/7, 4/8} \draw[ink2, line width=0.7pt] (e\e) -- (f\f);
  \node[lbl, anchor=east] at (-0.45,1.6) {flights};
  \node[lbl, anchor=east] at (-0.45,0) {episodes};
  \draw[decorate, decoration={brace, amplitude=5pt}, blueI] (-0.2,2.05) -- (3.2,2.05) node[midway, above=5pt, lbl, ink] {block $\{1,2,3,4\}$};
  \node[lbl, above=11pt] at (f5) {block $\{5\}$};
  \draw[decorate, decoration={brace, amplitude=5pt}, orangeJ] (4.8,2.05) -- (7.2,2.05) node[midway, above=5pt, lbl, ink] {block $\{6,7,8\}$};
\end{tikzpicture}
\caption{Dependency blocks of a day. Flight 3 is in episodes E1 and E2, so flights 1--4 form one
block although 1 and 4 share no episode; flight 5 never flies in a group.}
\label{fig:blocks}
\end{figure}

\begin{proposition}
The blocks are the finest partition of the flights in which the members of every episode are in
the same block.
\end{proposition}

\begin{proof}
A partition with this property must contain the transitive closure of the relation ``share an
episode'', whose classes are exactly the connected components of the bipartite graph; these, in
turn, satisfy the property.
\end{proof}

The blocks are the units resampled in the bootstrap (Section~\ref{sec:comparison}): flights of
different blocks never shared an episode. Buchin et al.\ already note that a grouping structure
splits the trajectories into independent data sets~\cite[§1]{buchin2015}.

\subsection{An example day}\label{sec:example-day}

Figure~\ref{fig:day} follows a day with five flights. At 11:40 flights 1, 8 and 22 meet (edges
1--8 and 8--22): configuration $c_1=\{1,8,22\}$ starts. At 12:05 the edge 1--22 closes the
triangle: the members do not change and $c_1$ goes on. At 12:55 the interval 1--8 ends as 1
changes route: $c_1$ ends and $c_2=\{8,22\}$ starts; they share 8 and 22, so $c_1\to c_2$. At
13:30 the edge 8--14 opens: $c_3=\{8,14,22\}$, and $c_2\to c_3$. At 14:10 the group dissolves and
episode E1 ends. At 15:00 flights 8 and 22 meet again: $c_4=\{8,22\}$ has no predecessor and
opens episode E2.

So flights 1 and 14 share an episode without ever sharing a configuration; 8 and 22 belong to two
episodes; 1, 8, 14 and 22 form one dependency block, and flight 5, never in a group, a block of
its own.

What we keep from this section: at every instant each flight has one label, \st{L}, \st{S} or its
configuration; configurations chain into episodes, and episodes tie the flights into blocks.
These labels are the output of the method (Appendix~\ref{sec:outputs}). They still depend on
$\tau^*$, which Section~\ref{sec:calibration} fixes.

\begin{figure}
\centering
\animated{example_day}
\caption{\animmark An example day; the animation runs the day from 11:00 to 16:00 and stops at
each event. Top: the active graph $G(t)$ at the five dotted instants; filled nodes are in a
configuration, coloured by episode; dashed nodes are in the local phase, faint ones not in flight. Bottom: the membership $\gamma_i(t)$ of each flight:
hatched is the local phase \st{L}, grey is solo \st{S}, coloured segments are the configurations
$c_1$--$c_4$, and the last row spans the episodes.}
\label{fig:day}
\end{figure}
